\documentclass[11pt]{article}
\usepackage{cmap}
\usepackage[margin=1.1in]{geometry}
\usepackage{amsmath,amssymb,amsthm}
\usepackage[T1]{fontenc}
\usepackage{lmodern}
\input{glyphtounicode}
\pdfgentounicode=1
\usepackage{xcolor}
\usepackage{hyperref}
\definecolor{linkblue}{rgb}{0.1,0.2,0.55}
\hypersetup{colorlinks=true,linkcolor=linkblue,citecolor=linkblue,urlcolor=linkblue}
\newtheorem{theorem}{Theorem}
\newtheorem{corollary}{Corollary}
\newtheorem{definition}{Definition}
\theoremstyle{remark}
\newtheorem{remark}{Remark}
\newcommand{\Var}{\mathrm{Var}}

\title{Parity Barriers for Decoupling Inequalities:\\
\large why no comparison functional of bounded marginal order can
certify uniform decoupling}
\author{Lluis Eriksson}
\date{July 2026 --- v3}

\begin{document}
\maketitle

\begin{abstract}
For every $r\ge1$, the uniform measure on the even-parity subset of
$\{\pm1\}^{r+1}$ is $r$-wise independent, yet the last coordinate has
unit variance while being an a.s.\ function of the others. This
example is classical --- parity-check codes are the standard
construction of $k$-wise independent distributions in the
pseudorandomness literature \cite{Joffe, ABI, AGM, AlonSpencer} ---
and we claim no novelty for it. What we record here is a consequence
we have found useful and have not seen isolated as a statement: any
``comparison functional'' whose value depends only on marginal data
of order $\le r$, with constants uniform over finite measures, takes
identical values on the parity measure and on the uniform product
measure, and is therefore consistent with perfect decoupling on a
measure where decoupling fails maximally. Hence no inequality built
from bounded-order functionals can imply uniform decoupling
principles --- Dobrushin-type mixing, approximate tensorisation with
measure-free constants, covariance decay --- on any class of measures
containing the parity family. The case $r=1$ recovers, and explains
structurally, the failure of raw-oscillation/Doob and
Efron--Stein-type steps found repeatedly in an adversarial audit of a
constructive Yang--Mills programme: those inequalities lie in the
order-$1$ class, and no repair within bounded-order data can succeed,
because the barrier recurs at every order. The finite statements are
kernel-checked in Lean~4/Mathlib by exhaustive computation
(\texttt{decide}) for $r\le4$ and verified in exact rational
arithmetic for $r\le6$; the general proof is a one-line character
computation, included for completeness.
\end{abstract}

\section{The classical example, and what is taken from it}

Fix $r \ge 1$ and $n = r+1$. Let $u_n$ be the uniform measure on
$\{\pm1\}^n$ and let $\pi_n$ be the uniform measure on the even-parity
set $\{x \in \{\pm1\}^n : \prod_i x_i = 1\}$. Two facts:

\begin{itemize}
\item[(a)] every subset of at most $r$ coordinates of $\pi_n$ is
distributed exactly as under $u_n$ ($\pi_n$ is \emph{$r$-wise
independent});
\item[(b)] $f(x) = x_n$ satisfies $\Var_{\pi_n}(f) = 1$ and
$f = \prod_{i \le r} x_i$ $\pi_n$-almost surely.
\end{itemize}

Both are classical. The support of $\pi_n$ is the parity-check code;
distributions supported on codes of dual distance $> k$ are $k$-wise
independent, and the parity example is the standard witness that
$(n-1)$-wise independence does not imply independence
\cite{Joffe, ABI, AGM, AlonSpencer}. For completeness, the Fourier
proof: $d\pi_n/du_n = 1 + \chi_{[n]}$ where
$\chi_S(x) = \prod_{i\in S} x_i$; for $1 \le |S| \le r$, character
orthogonality gives $\mathbb{E}_{\pi_n}[\chi_S] =
\mathbb{E}_{u_n}[\chi_S] + \mathbb{E}_{u_n}[\chi_{[n]\setminus S}] =
0$, which is (a); (b) holds because $x_n = \prod_{i\le r} x_i$ on the
support and the one-dimensional marginal of $x_n$ is uniform by (a).

The purpose of this note is the \emph{barrier} that (a)+(b) impose on
a specific family of inequalities that arise in constructive quantum
field theory, which we now formalise.

\section{The barrier}

\begin{definition}[comparison functional of order $r$]
\label{def:order}
Let $\mathcal M_n$ be the finite probability measures on
$\prod_{i=1}^n E_i$ (finite alphabets $E_i$). A functional
$\Phi_b(\mu,f)$ at site $b$ \emph{has order $r$} if it depends on
$(\mu, f)$ only through the joint marginals of $\mu$ of order $\le r$
together with the single-site data of $f$ at $b$ (oscillation,
conditional variance given order-$\le r$ marginal information). An
\emph{order-$r$ comparison inequality} is a bound of the form
$\mathcal D(\mu,f) \le c \sum_b \Phi_b(\mu,f)$ with $c$ uniform over
$\mathcal M_n$, where $\mathcal D$ is a decoupling functional
(variance defect, covariance, martingale increment).
\end{definition}

\begin{theorem}[parity barrier]
\label{thm:main}
Fix $r \ge 1$. No order-$r$ comparison inequality can bound any
decoupling functional that separates $\pi_{r+1}$ from $u_{r+1}$; in
particular none implies Dobrushin-type mixing, approximate
tensorisation with measure-independent constants, or covariance
decay, on any class of measures containing
$\{\pi_{r+1}\} \cup \{u_{r+1}\}$.
\end{theorem}

\begin{proof}
By (a), every order-$r$ functional takes identical values on
$(\pi_{r+1}, f)$ and $(u_{r+1}, f)$. Under $u_{r+1}$ the coordinate
$f(x) = x_{r+1}$ is independent of the others and every decoupling
functional is consistent with perfect decoupling (covariances vanish,
tensorisation is exact); under $\pi_{r+1}$ decoupling fails maximally
by (b). A uniform constant $c$ would have to certify both
simultaneously --- impossible for any $\mathcal D$ separating the
two, e.g.\ $\mathcal D = \Var(f) - \sum_b \mathbb E\,
\Var_b(f \mid \text{rest})$, which is $0$ under $u_{r+1}$ and $1$
under $\pi_{r+1}$.
\end{proof}

The construction generalises to arbitrary finite abelian alphabets
(take density $1 + \mathrm{Re}\,\chi$ for any full-support character
$\chi$: the same barrier at order $n-1$), and vector-valued
observables change nothing (apply the theorem coordinatewise).

\begin{corollary}[the audit failures, explained]
\label{cor:audit}
The raw-oscillation/Doob increment bound and the Efron--Stein
covariance step --- found, and refuted by the $r=1$ case, in four
papers of the 2602-series adversarial audit of a constructive
Yang--Mills programme \cite{series} --- are order-$1$ comparison
inequalities; the ``XOR'' counterexample used there is the $r=2$
barrier. No repair that stays within bounded-order data ---
including correlation-decorated variants using pair moments --- can
succeed: the barrier recurs at every order.
\end{corollary}

\begin{remark}[sharpness: what the barrier does \emph{not} forbid]
\label{rem:scope}
Inequalities whose constants consume the \emph{conditional} structure
of the measure (Dobrushin interdependence coefficients, log-Sobolev
constants of the true single-site conditionals) are of unbounded
order in the sense of Definition~\ref{def:order}, and are untouched:
indeed the classical positive results (Dobrushin uniqueness
$\Rightarrow$ tensorisation with constant $(1-\delta)^{-1}$
\cite{Dob}) are of exactly this type. The moral is a division of
labour: bounded-order data can \emph{falsify} decoupling, never
\emph{certify} it; certification must consume conditional
information. We believe this is the precise structural reason why
the decoupling step is the irreducible analytic core of constructive
Yang--Mills programmes, and why it cannot be bypassed by moment or
oscillation bookkeeping of any fixed order.
\end{remark}

\section{Machine verification}

Statements (a) and (b) are finite for each $r$ and hence decidable by
exhaustive computation. They are kernel-checked in Lean~4/Mathlib via
\texttt{decide} for $r = 1, 2, 3, 4$ (file
\texttt{ParityBarrier.lean} accompanying this note: theorems
\texttt{parity\_kwise\_r1}--\texttt{r4} for the vanishing of all
character sums of nonempty support and order $\le r$ over the parity
set, \texttt{parity\_support\_r1}--\texttt{r4} for the support
identity, and mean-zero statements giving the unit variance), and
independently verified in exact rational arithmetic for $r \le 6$
(script \texttt{verify\_parity.py}). The Lean file was checked with
\texttt{lake env lean} against Lean \texttt{4.30.0-rc2}, Mathlib
commit \texttt{cd3b69b}; the axiom oracle returns Lean's three
standard axioms for every theorem in the file, e.g.
\begin{verbatim}
'parity_kwise_r4' depends on axioms:
    [propext, Classical.choice, Quot.sound]
'parity_support_r4' depends on axioms:
    [propext, Classical.choice, Quot.sound]
'parity_mean_zero_r4' depends on axioms:
    [propext, Classical.choice, Quot.sound]
\end{verbatim}
(exit code 0). The general-$r$ statement has
the one-line proof above; a parametric-in-$r$ Lean proof (induction
over character orthogonality) is left as future work.

\subsection*{Acknowledgements}
This note is part of an audit-first research programme whose records,
Lean sources and numerical suites are public at
\url{https://github.com/lluiseriksson/THE-ERIKSSON-PROGRAMME}. The
parity example itself is classical and no priority is claimed for it;
the debt to the $k$-wise independence literature is recorded in the
references.

\begin{thebibliography}{9}\small

\bibitem{Joffe}
A.~Joffe,
\emph{On a set of almost deterministic $k$-independent random
variables},
Ann. Probab. \textbf{2} (1974), 161--162.

\bibitem{ABI}
N.~Alon, L.~Babai and A.~Itai,
\emph{A fast and simple randomized parallel algorithm for the maximal
independent set problem},
J. Algorithms \textbf{7} (1986), 567--583.

\bibitem{AGM}
N.~Alon, O.~Goldreich and Y.~Mansour,
\emph{Almost $k$-wise independence versus $k$-wise independence},
Inform. Process. Lett. \textbf{88} (2003), 107--110.

\bibitem{AlonSpencer}
N.~Alon and J.~H. Spencer,
\emph{The Probabilistic Method},
4th ed., Wiley, 2016.

\bibitem{Dob}
R.~L. Dobrushin,
\emph{The description of a random field by means of conditional
probabilities and conditions of its regularity},
Theory Probab. Appl. \textbf{13} (1968), 197--224.

\bibitem{ES}
B.~Efron and C.~Stein,
\emph{The jackknife estimate of variance},
Ann. Statist. \textbf{9} (1981), 586--596.

\bibitem{series}
L.~Eriksson,
\emph{THE MASTER MAP: an audit-first navigation guide to the
conditional construction of 4D $SU(N)$ Yang--Mills with mass gap},
ai.viXra:2602.0096, and the audited 2602-series;
repository \url{https://github.com/lluiseriksson/THE-ERIKSSON-PROGRAMME}.

\end{thebibliography}

\end{document}
